%% kapa: exam=1 round=31 year=2020 period=1 form=A subject=civil_law subject_ko=민법 seq=24 number=24
%% source: Q-Net 감정평가사 기출문제 (cst00309) artlSeq=5207219 file=extracted/1008/5207219/A형/A형/2020년도 제31회 감정평가사 1차 시험_1교시_A형.pdf
%% conversion: kordoc 4.24.51 -> tools/gen1.py (auto); review=proofread-sonnet; flags=none
\begin{question}{24}
등기에 관한 설명으로 옳지 않은 것은? (다툼이 있으면 판례에 따름)
\begin{choices}
  \item 경정등기는 원시적으로 진실한 권리관계와 등기가 일부 어긋나는 경우 이를 바로잡는 등기이다.
  \item 소유자만이 진정명의회복을 위한 소유권이전등기를 청구할 수 있다.
  \item 진정명의회복을 위한 소유권이전등기청구의 상대방은 현재의 등기명의인이다.
  \item 증여로 부동산을 취득하였으나 등기원인을 매매로 기재하였다면 그 등기는 무효등기이다.
  \item 그 이유가 무엇이든 당사자가 자발적으로 말소등기한 경우 말소회복등기를 할 수 없다.
\end{choices}
\end{question}
\answer{4}
